\section{Discussion and Limitations}
\paragraph{Using the bounds.} The bounds split any map's error into a part the tree forces and a part its layout chooses. Order error above $\nu^*$ can always be removed without changing the tree, since midpoint ordering attains $\nu^*$; the width-aware and bounded-topology bounds need not be attained, so an excess above them may lie partly in the bound. Because they need only the tree, the projected locations and the widths, the bounds can be computed before a map is drawn, to choose a direction or to decide whether a budget is worth spending, and any new layout of this family can be measured against them.

\paragraph{Position and neighborhoods.} A fixed spatial reference and a geometry-oriented embedding answer different questions. AnchorMTM supports reading positions and motion along a chosen direction; ST-MTM's higher reported TW shows the value of preserving local neighborhoods, and neither property guarantees the other. Our results support this task-specific claim, not universal superiority in spatial or temporal fidelity. One axis also carries one coordinate: reading two-dimensional positions would need a second map along another direction, with its own bound.

\paragraph{Limitations.} Author-processed fields and their small-sample TW protocol are not fully available, so published SNS and TW cannot support an identical-input significance test, and the two suites differ in preprocessing. The reference and the tracks are computed offline from the whole sequence, and support overlap is a correspondence proxy, not a physical identity. Order search is heuristic, and a continuous certificate does not remove raster rounding, so we check the rendered merge matrices but not every augmented internal label. The motion reading assumes known identities. We have not measured perceptual readability, the accuracy of human motion judgments or generalization to unseen data; these need user studies.

\section{Conclusion}
Anchoring the axis of a merge tree map to space has an unavoidable positional price. Writing merge tree maps as one layout problem, with TMTM and ST-MTM as exact-topology instances, we bound this price for every instance: the bounds are attained on 606 of 614 planar frames, and 31--95\% of the order error of TMTM and ST-MTM is not forced by the tree. AnchorMTM pays little more than the price: at the same zero merge error it lowers the mean worst displacement by 73--96\% on the ST-MTM benchmark, and a declared merge budget, verified on every rendered frame, buys further position where conflicts are shallow, while ST-MTM's published neighborhood preservation remains higher.
